\documentclass[../../main.tex]{subfiles}
\graphicspath{{../../image/}{../../figure/}}
\begin{document}

\section{波追踪}\label{section:ch13-波追踪}

本节旨在追踪近似解 $U_h$ 的各个波，并监测其强度与传播速度的演化。这并非易事：波相互作用既可能使同一特征族中相碰的波融合或消亡，也可能孕育出其他特征族的新波。

为获得直观，考察简单 Burgers 方程
$\partial_t u+\frac12\partial_x u^2=0$%
\cref{chapter:ch04-cauchy问题}，

菱形中的波相互作用。此时波相互作用估计\eqref{align:ch13-int-main}退化为
\begin{align}
\varepsilon=\alpha+\beta.
\end{align}

在一种典型情形中，(负)振幅为 $\alpha,\beta$、速度分别为 $\sigma_L,\sigma_R$ 的激波，分别通过菱形的``西南''边与``东南''边进入，并融合成一个振幅为 $\varepsilon=\alpha+\beta$、速度为 $\sigma$ 的激波。有启发性的做法是把出射激波视为两个``虚拟波''的复合，其振幅分别为 $\alpha$ 与 $\beta$，从而两入射激波在碰撞后仍继续前行，振幅不变而速度改变。由 $\sigma\varepsilon=\sigma_L\alpha+\sigma_R\beta$，易得
\begin{align}
|\sigma-\sigma_L||\alpha|=|\sigma-\sigma_R||\beta|=\frac12\alpha\beta.
\label{align:ch13-wt-shockcancel}
\end{align}
回忆 $\alpha\beta$ 即该菱形中的波相互作用量。

在对偶情形中，(正)振幅为 $\alpha,\beta$ 的稀疏波分别通过菱形的``西南''边与``东南''边进入，合并成一个振幅为 $\varepsilon=\alpha+\beta$ 的稀疏波，后者又分裂为振幅为 $\alpha',\beta'$ 的新稀疏波，分别通过菱形的``西北''边与``东北''边离开。例如设 $\alpha'<\alpha$，则可把左入射波视为两个稀疏波的复合，振幅分别为 $\alpha'$ 与 $\alpha-\alpha'$；而右出射波视为两个稀疏波的复合，振幅分别为 $\alpha-\alpha'$ 与 $\beta$。如此，三入射波在相互作用后均以不变振幅、改变了的速度继续前行。

还有一种情形：(负)振幅为 $\alpha$、速度为 $\sigma_L$ 的激波通过``西南''边进入菱形，与通过``东南''边进入的(正)振幅为 $\beta$ 的稀疏波相互作用。例如设 $|\alpha|>|\beta|$，则出射波是振幅为 $\varepsilon=\alpha+\beta$、速度为 $\sigma$ 的激波。如前，把入射激波视为两个``虚拟波''的复合，振幅分别为 $\alpha+\beta$ 与 $-\beta$。于是相互作用的结果是：第二个入射虚拟波与入射稀疏波相互对消，而第一个虚拟波以不变振幅、改变了的速度继续前行。简单计算给出速度改变量
\begin{align}
|\sigma-\sigma_L|=\frac12\beta.
\label{align:ch13-wt-speedchange}
\end{align}
注意 $\beta$ 即该菱形中的波对消量。

离开上述菱形的波将卷入后续碰撞，并在其中可能须进一步分割为更细的虚拟波。这些分割应沿时间反向追溯，施加到该波的每一个祖先上。这一繁复过程的最终结果是：在任一指定时间区域内，每个波都被分割为若干虚拟波，它们归为以下两类之一：在指定时间区间内历经所有碰撞后存留者，以及最终经相互对消而湮灭者。

对双曲守恒律系统，情形类似，但须注意：任意两波的碰撞一般会孕育出每一特征族的新波。在特征族分段真非线性或线性退化的严格双曲系统中，波按如下程序分割为虚拟波。

连接左状态 $U_-$ 与右状态 $U_+$ 的 $i$ 激波的一个分割，由某状态序列 $U_-=U^0,U^1,\cdots,U^\nu=U_+$ 实现，其中 $\mu=1,\cdots,\nu$，$U^\mu$ 位于从 $U_-$ 发出的 $i$ 激波曲线上，且 $\lambda_i(U^\mu)\leqslant\lambda_i(U^{\mu-1})$。尽管 $U^{\mu-1}$ 与 $U^\mu$ 一般并不由激波连接，仍把数对 $(U^{\mu-1},U^\mu)$ 视为一个虚拟波，其振幅为 $V_i^\mu=U^\mu-U^{\mu-1}$，速度为 $\lambda_i^\mu$，等于激波 $(U_-,U_+)$ 的速度。

连接左状态 $U_-$ 与右状态 $U_+$ 的 $i$ 稀疏波的分割类似地由有限状态序列实现，即 $U_-=U^0,U^1,\cdots,U^\nu=U_+$，其中 $\mu=1,\cdots,\nu$，$U^\mu$ 位于从 $U_-$ 发出的 $i$ 稀疏波曲线上，且 $\lambda_i(U^\mu)>\lambda_i(U^{\mu-1})$。尽管此时 $U^{\mu-1}$ 与 $U^\mu$ 可由真实的 $i$ 稀疏波连接，$(U^{\mu-1},U^\mu)$ 仍视为虚拟波，振幅为 $V_i^\mu=U^\mu-U^{\mu-1}$，速度为 $\lambda_i^\mu=\lambda_i(U^{\mu-1})$。

连接左状态 $U_-$ 与右状态 $U_+$ 的一般 $i$ 波，由有限个 $i$ 激波与 $i$ 稀疏波组成，其分割以显然方式组合上述纯激波与纯稀疏波两种情形而得到。

借助\cref{section:ch13-注记}所引文献中的繁复构造，在指定时间区域
\begin{align*}
\Lambda=\{(x,t):-\infty<x<\infty,\ s_1\lambda^{-1}h\leqslant t\leqslant s_2\lambda^{-1}h\}
\end{align*}
上，近似解 $U_h$ 的波可分割为分属以下三类之一的虚拟波：

\noindent\textbf{I.}\, 于 $t=s_1\lambda^{-1}h$ 以正强度进入 $\Lambda$，在时间区间 $[s_1\lambda^{-1}h,s_2\lambda^{-1}h]$ 内存活，并于 $t=s_2\lambda^{-1}h$ 以正强度离开 $\Lambda$ 的波。

\noindent\textbf{II.}\, 于 $t=s_1\lambda^{-1}h$ 以正强度进入 $\Lambda$，但在 $\Lambda$ 内经相互对消而湮灭的波。

\noindent\textbf{III.}\, 在 $\Lambda$ 内部通过波相互作用产生的波。

设 $\mathscr{W}$ 表示上述三类中任一类的典型虚拟波，目标是估计其在 $\Lambda$ 内存活期内的最大强度，记作 $|\mathscr{W}|$；振幅的总变差，记作 $[\mathscr{W}]$；以及速度的总变差，记作 $[\sigma(\mathscr{W})]$。这些估计的依据在于标量情形得到的简单估计\eqref{align:ch13-wt-shockcancel}与\eqref{align:ch13-wt-speedchange}，结合前几节导出的波相互作用估计。

对特征族真非线性的系统，所需估计为
\begin{align}
\sum_{\mathscr{W}\in\mathrm{I}}\big\{[\mathscr{W}]+|\mathscr{W}|[\sigma(\mathscr{W})]\big\}=O(1)\mathscr{D}(\Lambda),
\label{align:ch13-wt-classI}\\
\sum_{\mathscr{W}\in\mathrm{II}}\big\{[\mathscr{W}]+|\mathscr{W}|\big\}=O(1)\mathscr{C}(\Lambda)+O(1)\mathscr{D}(\Lambda),
\label{align:ch13-wt-classII}\\
\sum_{\mathscr{W}\in\mathrm{III}}\big\{[\mathscr{W}]+|\mathscr{W}|\big\}=O(1)\mathscr{D}(\Lambda),
\label{align:ch13-wt-classIII}
\end{align}
其中 $\mathscr{D}(\Lambda)$ 与 $\mathscr{C}(\Lambda)$ 分别表示 $\Lambda$ 内波相互作用与波对消的总量，即
\begin{align}
\mathscr{D}(\Lambda)=\sum\mathscr{D}(\Delta_s^r),\qquad
\mathscr{C}(\Lambda)=\sum\mathscr{C}(\Delta_s^r),
\label{align:ch13-wt-DC}
\end{align}
求和遍历所有包含于 $\Lambda$ 内的菱形 $\Delta_s^r$，而 $\mathscr{D}(\Delta_s^r)$、$\mathscr{C}(\Delta_s^r)$ 由\eqref{align:ch13-D-def}、\eqref{align:ch13-C-def}定义。

对特征族仅为分段真非线性的系统，\eqref{align:ch13-wt-classI}、\eqref{align:ch13-wt-classII}与\eqref{align:ch13-wt-classIII}的类似估计要复杂得多。差异源于：波相互作用量 $\mathscr{D}(\Delta_s^r)$ 在真非线性情形为二阶量\eqref{align:ch13-D-def}，而在一般情形仅为三阶量\eqref{align:ch13-gen-D}。细节见\cref{section:ch13-注记}所引文献。

由此可建立以下命题，它改进了\cref{theorem:ch13-几乎处处相容}，去掉了序列 $\wp$ 选取中的``随机性''假设：

\begin{theorem}[等分布序列的相容性]\label{theorem:ch13-等分布相容}
由任一等分布于区间 $(-1,1)$（在\eqref{align:ch13-equidist}意义下）的序列 $\wp=\{a_0,a_1,\cdots\}$ 所诱导的算法是相容的。
\end{theorem}

\begin{proof}
证明可在\cref{section:ch13-注记}所引文献中找到：将\eqref{align:ch13-consistency}右端用分割 $U_h$ 的虚拟波表出，进而证明只要序列 $\wp$ 等分布，当 $h\downarrow0$ 时它趋于零。其理由如下。回忆在\cref{section:ch13-紧性与相容性}中，我们已对线性守恒律 $\partial_t u+a\partial_x u=0$，利用每个波以常振幅、常速度传播的性质，验证了算法对任一等分布序列 $\wp$ 的相容性。上述波分割表明，即使非线性系统也近似具有这一性质，从而可把相容性论证推广到该情形。证毕。
\end{proof}

波划分虽略显繁复，却是获得解的局部结构、大时间行为及其他定性性质精确信息的有效工具；尤其对于依赖全局波型的性质的推导，它不可或缺。

\end{document}
